\documentclass{beamer}
\usetheme{ut}

% --------------------------------------------------
% Encoding / language
% --------------------------------------------------
\usepackage[T1]{fontenc}
\usepackage{lmodern}
\usepackage[utf8]{inputenc}
\usepackage[english]{babel}
\usepackage{csquotes}
\usepackage{amsmath,amssymb}
\usepackage{graphicx}
\usepackage{booktabs}
\usepackage{array}
\usepackage{tikz}
\usepackage{hyperref}

% --------------------------------------------------
% Metadata
% --------------------------------------------------
\title[ Topology-resolved quantum volume ]{ From a single number to a landscape: topology-resolved quantum volume }
\subtitle[Quantum Science and Technologies]{Quantum Science and Technologies}
\author[H.~Kravals, M.~Dimitrijevs]{Henrijs Kravals, Maksims Dimitrijevs}
\institute[LU]{Latvijas Universitāte}
\date[Apr 2026]{April 15, 2026}
\footlinetext{[ Topology-resolved quantum volume – H.~Kravals, M.~Dimitrijevs \;]}

\begin{document}

\maketitle

% ==================================================
% CHAPTER 1
% ==================================================

\begin{frame}{Goal of the presentation}
    \begin{columns}[T,totalwidth=\textwidth]
        \column{0.4\textwidth}
        

        \column{0.4\textwidth}
        \begin{block}{Structure}
            \begin{enumerate}
                \item NISQ computers,
                \item Noisy simulators,
                \item Quantum volume,
                \item Research demo
            \end{enumerate}
        \end{block}
    \end{columns}
\end{frame}


\begin{frame}{NISQ computers}

    

\end{frame}

\begin{frame}{Noisy simulators}

    

\end{frame}


\begin{frame}{Quantum volume (1)}
    Quantum volume (QV) is a \textbf{system level benchmark}.
    
    \begin{itemize}
        \item<2-> The QV is a single number figure of merit used to evaluate universal gate-based quantum computers. 
    \end{itemize}

    \only<3->{
        \begin{exampleblock}{\emph{Heavy set of outputs}}
            Let $\mathbf{Q}$ be a random circuit drawn from a suitable ensemble acting on $n$ qubits. The quantum state after executing the circuit is denoted $|\psi\rangle$. Each possible output state $x\in\{0,1\}^n$ is measured with probability $|\langle x|\psi\rangle|^2$. The set of output states with a probability greater than the median constitutes the heavy set of outputs associated with the quantum circuit $\mathbf{Q}$.
        \end{exampleblock}
    }
    
\end{frame}


\begin{frame}{Quantum volume (2)}
    \begin{figure}
        \centering
        \includegraphics[width=\linewidth,height=0.8\textheight,keepaspectratio]{images/qv-process.png}
        \caption{The QV circuit. Source: \url{https://quantumbenchmarkzoo.org/content/system-level-benchmark/others/quantum-volume}}
    \end{figure}
\end{frame}


\begin{frame}{Quantum volume (3)}
    If the quantum computer samples a distribution from the \textbf{heavy set of outputs} (simulated classically) with probability $> 2/3$, it validates the corresponding quantum volume score of $2^n$. The confidence interval for this evaluation is set to $2\sigma$ ($97.7\%$).

    \begin{alertblock}{Problem}
        Quantum volume assesses only a limited subset of best-quality qubits within a quantum processor and does not provide a comprehensive measure of the overall fidelity or performance of quantum operations across the entire chip.
    \end{alertblock}
\end{frame}


\begin{frame}{Research methodology}
    We benchmarked QV on the IBM Marrakesh device using calibration data and the coupling graph from Qiskit.

    \begin{enumerate}
        \item Construct a noisy quantum simulator from the extracted data
        \item For each size $n \leq 6$:
        \begin{enumerate}
            \item Enumerate all connected subgraphs of size $n$
            \item Run QV circuits on each subgraph
            \item Compute heavy output probability and confidence
        \end{enumerate}
    \end{enumerate}
\end{frame}

\begin{frame}{Results}
    \small
    \begin{tabular}{rrlllllll}
\toprule
$n$ & Count & $\langle HOP \rangle$ & $\sigma(HOP)$ & Median & Min & Max & $\langle SE \rangle$ & Lower \\
\midrule
3 & 244 & 0.664 & 0.063 & 0.688 & 0.438 & 0.712 & 0.047 & 0.572 \\
4 & 354 & 0.732 & 0.083 & 0.766 & 0.503 & 0.817 & 0.043 & 0.647 \\
5 & 566 & 0.733 & 0.085 & 0.769 & 0.498 & 0.825 & 0.043 & 0.648 \\
6 & 926 & 0.748 & 0.085 & 0.779 & 0.491 & 0.832 & 0.042 & 0.665 \\
\bottomrule
\end{tabular}

\end{frame}

\begin{frame}{Discussion}
    
For QV, each layer is built from random SU(4) two-qubit blocks after a random permutation, and for odd widths one qubit is idle each layer. On a connected 3-qubit path like (a, b, c), one of the possible two-qubit interactions is between the two endpoints. That interaction is not native, so with trivial layout and basic routing it tends to introduce SWAP overheadresulting in extra routing cost that can dominate the whole layer. 3Q layers contain only one gate, so routing dominates. 4Q layers contain two gates, allowing routing overhead to be shared.

    

\end{frame}


\maketitle

\end{document}